\documentclass[a4paper]{article}

% 文章排版
\usepackage{indentfirst}
\setlength{\parindent}{0em}
\usepackage{autobreak}
\allowdisplaybreaks
\usepackage{parskip}
\usepackage{geometry}
\geometry{top=3cm,bottom=3cm,left=2cm,right=2cm}
\usepackage{anyfontsize}

% 数学物理宏包
\usepackage{amsmath}
\usepackage{physics}
\renewcommand\div\divisionsymbol
\DeclareMathOperator{\diag}{diag}
\DeclareMathOperator{\sgn}{sgn}
\newcommand{\trans}{\text{\textsf{T}}}
\usepackage{tabularray}
\UseTblrLibrary{amsmath}

\usepackage{xcolor}
\usepackage{fontspec}
\setmainfont{STIX Two Text}
\usepackage{unicode-math}
\unimathsetup{mathrm=sym,mathbf=sym,mathsf=sym,bold-style=ISO}
\setmathfont{STIX Two Math}

\usepackage[fontset=none]{ctex}
\setCJKmainfont{FZShuSong-Z01}[BoldFont=Source Han Sans CN,ItalicFont=FZKai-Z03]
\setCJKsansfont{Source Han Sans CN}

% 符号重定义
\AtBeginDocument{
    \let\uglyepsilon\epsilon
    \renewcommand\epsilon\varepsilon
    \let\uglyleq\leq
    \renewcommand\leq\leqslant
    \let\uglygeq\geq
    \renewcommand\geq\geqslant
}


% 题目
\title{\textbf{广义相对论（天文方向）\ \ \ \ 作业七}}
\author{国家天文台 张植竣}
\date{2025年10月24日}

\begin{document}

\maketitle

\begin{enumerate}
    \item 证明Einstein场方程：
    \[
        R_{\mu\nu} - \frac{1}{2} g_{\mu\nu} R = \kappa T_{\mu\nu}
    \]
    可以写为：
    \[
        R_{\mu\nu} = \kappa \qty(T_{\mu\nu} - \frac{1}{2} g_{\mu\nu} T)
    \]
    其中$T = T^{\mu}_{\ \mu}$是能动张量的迹。

    \textcolor{blue}{
    比较两个等式：
    \[
        \begin{+cases}
            & R_{\mu\nu} = \dfrac{1}{2} g_{\mu\nu} R + \kappa T_{\mu\nu} \\
            & R_{\mu\nu} = -\dfrac{1}{2} g_{\mu\nu} \kappa T + \kappa T_{\mu\nu}
        \end{+cases}
    \]
    那么只需要证明$R = -\kappa T$即可。 \\
    证明如下，对Einstein场方程等号两边都作用$g^{\mu\nu}$：
    \[
        g^{\mu\nu} \qty(R_{\mu\nu} - \frac{1}{2} g_{\mu\nu} R) = g^{\mu\nu} \qty(\kappa T_{\mu\nu})
    \]
    由于：
    \[
        g^{\mu\nu} R_{\mu\nu} = R,\quad g^{\mu\nu} g_{\mu\nu} = \delta^\mu_{\ \mu} = 4,\quad g^{\mu\nu} T_{\mu\nu} = T
    \]
    该等式变为：
    \[
        R - \frac{1}{2} \times 4 R = -R = \kappa T
    \]
    因此：
    \[
        R_{\mu\nu} = \dfrac{1}{2} g_{\mu\nu} R + \kappa T_{\mu\nu} = -\dfrac{1}{2} g_{\mu\nu} \kappa T + \kappa T_{\mu\nu} = \kappa \qty(T_{\mu\nu} - \frac{1}{2} g_{\mu\nu} T)
    \]
    }

    \item 证明真空Einstein场方程可以写为$R_{\mu\nu} = 0$。

    \textcolor{blue}{
    真空Einstein场方程为：
    \[
        G_{\mu\nu} = R_{\mu\nu} - \frac{1}{2} g_{\mu\nu} R = 0
    \]
    对真空Einstein场方程等号两边都作用$g^{\mu\nu}$：
    \[
        g^{\mu\nu} \qty(R_{\mu\nu} - \frac{1}{2} g_{\mu\nu} R) = 0
    \]
    由于：
    \[
        g^{\mu\nu} R_{\mu\nu} = R,\quad g^{\mu\nu} g_{\mu\nu} = \delta^\mu_{\ \mu} = 4
    \]
    该等式变为：
    \[
        R - \frac{1}{2} \times 4 R = -R = 0
    \]
    即$R = 0$，因此：
    \[
        R_{\mu\nu} - \frac{1}{2} g_{\mu\nu} R = R_{\mu\nu} - \frac{1}{2} g_{\mu\nu} \times 0 = 0
    \]
    故真空Einstein场方程可以写为$R_{\mu\nu} = 0$。
    }

\end{enumerate}

\end{document}